The preceding sections developed the firm's problem over both horizons.
This section collects two further results that the framework delivers but that the empirical work of Section~\ref{sec:empirics} does not draw on.
Section~\ref{sec:extensions.aggregation} asks whether a technology built on sequential workflows can be aggregated into a production function for the economy, and shows that it can.
Section~\ref{sec:extensions.computation} asks whether the cost minimization problems we have posed can actually be solved, and shows that they can be, in polynomial time.


\subsection{CES Representation}
\label{sec:extensions.aggregation}

Everything the model has said so far happens inside a single firm: which steps it chains together, where it draws the boundaries of jobs, and what those jobs pay.
One might reasonably object that this is not where the effects of AI will be settled.
The consequences people care about are macroeconomic, namely how much labor of each kind the economy employs, how output responds to capital and to labor, and what becomes of both as AI improves, and that is where the automation literature works, writing production functions directly at the level of the economy \citep{acemoglu2018, acemoglu2022artificial, acemoglu2024tasks}.
A framework built on the internal reorganization of workflows might seem to have nothing to say in those terms, or worse, to sit awkwardly beside them.

It does not.
We show in Appendix~\ref{sec:aggregation} that the firm-level technology developed above, in which the boundaries of tasks and of jobs are chosen rather than given, aggregates to
\begin{align}
\label{eq:lr_macro_ces}
X = \left(\theta_{\AIletter}\,\aggregateAIlabor^{\rho} + \theta_{\manualLetter}\,\aggregateManualLabor^{\rho} + (1-\theta_{\AIletter}-\theta_{\manualLetter})\,K^{\rho}\right)^{\frac{1}{\rho}},
\end{align}
over economy-wide AI management labor $\aggregateAIlabor$, manual labor $\aggregateManualLabor$, and capital $K$.\footnote{
The representation asks more of the economy than the firm-level argument does: capital productivity is common across firms, every firm organizes identically, and realized effective AI quality is the single dimension along which firms differ.
}
This is the familiar form, the one in which the macroeconomic implications of automation are usually studied, with an elasticity of substitution between AI and human labor that carries most of the weight in any comparative static.

Arriving there is not automatic, and what makes it work is worth stating.
No firm in our model substitutes AI management labor for manual labor, since producing a unit of output requires completing every step of the workflow, so each firm's technology is Leontief.
Aggregate substitution comes instead from composition, in the manner of \citet{houthakker1955pareto} and \citet{levhari1968note}: firms have access to the same AI but differ in how effectively they deploy it, and when AI management labor becomes relatively cheaper the economy's input mix moves because the set of firms for which using AI is worthwhile has moved, not because any firm has changed its own mix.
Appendix~\ref{sec:aggregation} obtains that dispersion from the order in which a firm has to commit, settling on its AI strategy and job design before it learns how effectively it can put AI to use, and derives in closed form the distribution of realized AI effectiveness under which \eqref{eq:lr_macro_ces} holds.

The two levels are therefore not in competition.
Our account of how chaining redraws tasks and jobs is not an alternative to the aggregate view of automation but a foundation for it, and the implications people draw from improving AI quality in a CES economy carry over intact.
The organizational content carries over as well, but only in the parameters of \eqref{eq:lr_macro_ces}: two economies whose steps are identically exposed to AI, differing only in how those steps are arranged within workflows, reach the same functional form with different weights and input requirements while looking alike on any step-level exposure measure.
That is why workflow-level objects such as the fragmentation index have to be measured where the arrangement of work is still visible, which is what Section~\ref{sec:empirics} does.


\subsection{Computing the Firm's Optimum}
\label{sec:extensions.computation}

Throughout the paper we have taken the firm to be at its cost-minimizing arrangement without asking whether that arrangement can be found.
The question is not idle here, because chaining is what makes it hard.
Whether automating a step lowers cost depends on the arrangement the step sits in, since a step costs nothing directly once it is absorbed into a chain but adds its failure risk to every attempt at that chain, so the firm chooses among whole arrangements rather than deciding separately at each step, and the number of arrangements grows exponentially in the number of steps $m$.
A model in which the optimum is out of reach would be a poor positive account of what firms do, and it would also leave the fragmentation index of Section~\ref{sec:results} approximating a quantity no one could compute.
This subsection shows that neither concern applies: the short-run problem is solvable exactly and cheaply, and the long-run problem is solvable to any desired accuracy.

\paragraph{The short run.}
What makes the short-run problem \eqref{eq:totalcost_shortterm} tractable is that it adds up time costs across tasks.
Once the firm has decided where a task ends, the cost of everything before it is a separate problem, so the firm can walk the workflow from left to right and face only one question at each step: whether to perform that step manually, or to end an AI chain there, and in the latter case how far back the chain reaches.

\begin{proposition}[Optimal AI Deployment]
\label{prop:shortrun_dp}
Given a workflow with $m$ steps, the AI strategy minimizing \eqref{eq:totalcost_shortterm} can be calculated in time $O(m^2)$ via dynamic programming.
\end{proposition}

Appendix~\ref{app:shortrun_proof} gives the proof.
Writing $C[k]$ for the minimum cost of completing the first $k$ steps alone, the two options above give a recursion for $C[k]$ in terms of $C[0], \dotsc, C[k-1]$; each of the $m$ values costs $O(m)$ to evaluate because the chain reaching back from step $k$ can start anywhere before it, which is where the $O(m^2)$ comes from.

\paragraph{The long run.}
The long-run problem \eqref{eq:totalcost_with_handoff} is harder on two counts.
Its objective is not a sum of task costs but a sum of products of two sums, one for the skill a job requires and one for the time it takes, so a task's contribution to cost depends on what the rest of its job already carries and the recursion above no longer applies.
The firm must also decide where jobs end, not only where chains end.
Exactness is what has to give, but only by an amount the firm can choose.

\begin{proposition}[Long-run AI Deployment and Job Design]
\label{prop:totalcost_optimization_dp}
Fix a workflow with $m$ steps whose skill, time, and hand-off costs lie in $[1/B, B]$ for some $B > 0$.
For any $\epsilon > 0$, a pair of AI strategy $\T$ and job design $\J$ whose long-run cost~\eqref{eq:totalcost_with_handoff} is within a factor $(1+\epsilon)$ of the minimum can be computed in time $O(m^2 \epsilon^{-2} \log^2(mB))$ via dynamic programming.
\end{proposition}

The program still walks the workflow from left to right, but it now carries the accumulated skill and time of the worker currently being assigned tasks, and at each step chooses among three options rather than two: execute the step manually, end an AI chain there, or close that worker's job and start a new one, paying the hand-off cost at that boundary.
Closing a job is the point at which the accumulated skill and time are multiplied into a wage bill, which is why the state has to carry both.
That state can take exponentially many values, so we discretize each coordinate to a geometric grid, which keeps the running time polynomial at the cost of an error that can be made arbitrarily small.
Appendix~\ref{app:jobdesign.longrun} gives the proof, building on the intermediate case of Appendix~\ref{app:mediumrun_dp}, in which wages adjust to the skill a job requires but job boundaries are still given.

Taken together, these results say that the firm's problem is not merely well posed but solvable, in the short run exactly and in the long run to any accuracy the firm cares to specify.
This matters for how the rest of the paper should be read.
The predictions we take to the data in Section~\ref{sec:empirics} below are predictions about an optimum, and they would be of limited interest if that optimum were something no firm could locate.
It also suggests a use for the framework beyond what we do here: paired with detailed measurement of the time and skill content of individual tasks, these algorithms could be run on real workflows to ask how a given advance in AI capability would redraw them.
